\ssr{ВВЕДЕНИЕ}

В последние годы наблюдается стремительный рост количества датасетов,
используемых в научных исследованиях, бизнес-аналитике и машинном обучении.
Эти наборы данных часто хранятся в разрозненных источниках, без унифицированного подхода к версионированию и поиску.
Отсутствие централизованного и структурированного хранилища затрудняет повторное использование и совместную работу.
Это порождает необходимость в разработке базы данных, ориентированной на хранение и управление датасетами.


Цель работы --- создание программного обеспечения для хранения, версионирования и поиска наборов данных (датасетов).


Для достижения поставленной цели требуется решить следующие задачи:
\begin{itemize}
    \item провести анализ предметной области и определить ключевые сущности и их взаимосвязи;
    \item сформулировать функциональные и нефункциональные требования к системе;
    \item спроектировать структуру базы данных, включая ограничения целостности и ролевую модель;
    \item реализовать физическую модель базы данных и необходимые механизмы доступа;
    \item провести исследование влияния различных типов индексов на производительность выполнения запросов.
\end{itemize}